% GENERATED FILE. Edit canonical lessons, then run npm run generate:books.
% canonical-source-hash: fnv1a64:0aab6b352b37aa7a
% canonical-lessons: ES-C620-solucion, ES-C620-solucionar, ES-C620-obstaculo, ES-C620-complicado, ES-C620-plantear

\chapter{Buscar soluciones}
\label{ch:es-buscar-soluciones}
\hlchaptermodality{620}

\begin{chapteropening}
\textbf{By the end of this chapter:} \emph{I can raise a problem, describe the difficulty and propose a solution.}

Complete the last lesson of chapter 620: I can raise a problem, describe the difficulty and propose a solution.
\end{chapteropening}

\section[la solución]{la solución — the solution}

\label{lesson:ES-C620-solucion}



\begin{quote}
Before the new one: say the Spanish for the prospect, then the Spanish for to predict.
\end{quote}

\subsection*{You'll want to know: la solución}
\textbf{la solución} — \textquotedblleft{}the solution\textquotedblright{}.

\begin{quote}
Tenemos que buscar una solución.
\end{quote}

\begin{quote}
\emph{We have to look for a solution.}
\end{quote}

\emph{No tiene solución} — it cannot be fixed.

\begin{cousinweb}
\emph{solución} is a learned borrowing from Latin \emph{solutio} 'a loosening, a solving', built on \emph{solvere} 'to loosen' (EtymDB).
\end{cousinweb}

\begin{grammarlens}[title={What you've built}]
One more word for telling, explaining and describing at B1.
\end{grammarlens}

\subsection*{Your turn}
\begin{itemize}
\raggedright
  \item \emph{Say it:} \emph{la solución}
  \item \emph{Say it:} \emph{la solución}, once more
  \item \emph{Say it:} a sentence of your own with \emph{la solución}
  \item \emph{Recall:} say the Spanish for the prospect, then the Spanish for to predict, then say \emph{la solución} again
\end{itemize}

\begin{tcolorbox}[breakable,colback=teal!4,colframe=teal!35!black,title={Before you move on}]
What does la solución mean? (\textquotedblleft{}The solution\textquotedblright{}.) Say it once more.
\end{tcolorbox}
\section[solucionar]{solucionar — to solve, to sort out}

\label{lesson:ES-C620-solucionar}



\begin{quote}
Before the new one: say the Spanish for to predict, then the Spanish for the solution.
\end{quote}

\subsection*{You'll want to know: solucionar}
\textbf{solucionar} — \textquotedblleft{}to solve, to sort out\textquotedblright{}.

\begin{quote}
Vamos a solucionar el problema.
\end{quote}

\begin{quote}
\emph{We are going to sort out the problem.}
\end{quote}

Regular -ar verb; close to \emph{resolver}.

\begin{cousinweb}
\emph{solucionar} is built on \emph{solución}, a learned borrowing from Latin \emph{solutio}, from \emph{solvere} 'to loosen' (EtymDB).
\end{cousinweb}

\begin{grammarlens}[title={What you've built}]
One more word for telling, explaining and describing at B1.
\end{grammarlens}

\subsection*{Your turn}
\begin{itemize}
\raggedright
  \item \emph{Say it:} \emph{solucionar}
  \item \emph{Say it:} \emph{solucionar}, once more
  \item \emph{Say it:} a sentence of your own with \emph{solucionar}
  \item \emph{Recall:} say the Spanish for to predict, then the Spanish for the solution, then say \emph{solucionar} again
\end{itemize}

\begin{tcolorbox}[breakable,colback=teal!4,colframe=teal!35!black,title={Before you move on}]
What does solucionar mean? (\textquotedblleft{}To solve, to sort out\textquotedblright{}.) Say it once more.
\end{tcolorbox}
\section[el obstáculo]{el obstáculo — the obstacle}

\label{lesson:ES-C620-obstaculo}



\begin{quote}
Before the new one: say the Spanish for the solution, then the Spanish for to solve.
\end{quote}

\subsection*{You'll want to know: el obstáculo}
\textbf{el obstáculo} — \textquotedblleft{}the obstacle\textquotedblright{}.

\begin{quote}
Hay muchos obstáculos en el camino.
\end{quote}

\begin{quote}
\emph{There are many obstacles in the way.}
\end{quote}

\emph{Superar un obstáculo} — to overcome an obstacle.

\begin{cousinweb}
\emph{obstáculo} is a learned borrowing from Latin \emph{obstaculum} (EtymDB), built on \emph{obstare} 'to stand in the way', from \emph{stare} 'to stand'.
\end{cousinweb}

\begin{grammarlens}[title={What you've built}]
One more word for telling, explaining and describing at B1.
\end{grammarlens}

\subsection*{Your turn}
\begin{itemize}
\raggedright
  \item \emph{Say it:} \emph{el obstáculo}
  \item \emph{Say it:} \emph{el obstáculo}, once more
  \item \emph{Say it:} a sentence of your own with \emph{el obstáculo}
  \item \emph{Recall:} say the Spanish for the solution, then the Spanish for to solve, then say \emph{el obstáculo} again
\end{itemize}

\begin{tcolorbox}[breakable,colback=teal!4,colframe=teal!35!black,title={Before you move on}]
What does el obstáculo mean? (\textquotedblleft{}The obstacle\textquotedblright{}.) Say it once more.
\end{tcolorbox}
\section[complicado]{complicado — complicated}

\label{lesson:ES-C620-complicado}



\begin{quote}
Before the new one: say the Spanish for to solve, then the Spanish for the obstacle.
\end{quote}

\subsection*{You'll want to know: complicado}
\textbf{complicado} — \textquotedblleft{}complicated\textquotedblright{}.

\begin{quote}
La situación es complicada.
\end{quote}

\begin{quote}
\emph{The situation is complicated.}
\end{quote}

The participle of \emph{complicar}.

\begin{cousinweb}
\emph{complicado} is the participle of \emph{complicar}, a learned borrowing from Latin \emph{complicare} 'to fold together', built on \emph{plicare} 'to fold' (EtymDB).
\end{cousinweb}

\begin{grammarlens}[title={What you've built}]
One more word for telling, explaining and describing at B1.
\end{grammarlens}

\subsection*{Your turn}
\begin{itemize}
\raggedright
  \item \emph{Say it:} \emph{complicado}
  \item \emph{Say it:} \emph{complicado}, once more
  \item \emph{Say it:} a sentence of your own with \emph{complicado}
  \item \emph{Recall:} say the Spanish for to solve, then the Spanish for the obstacle, then say \emph{complicado} again
\end{itemize}

\begin{tcolorbox}[breakable,colback=teal!4,colframe=teal!35!black,title={Before you move on}]
What does complicado mean? (\textquotedblleft{}Complicated\textquotedblright{}.) Say it once more.
\end{tcolorbox}
\section[plantear]{plantear — to raise, to bring up; to pose}

\label{lesson:ES-C620-plantear}



\begin{quote}
Before the new one: say the Spanish for the obstacle, then the Spanish for complicated.
\end{quote}

\subsection*{You'll want to know: plantear}
\textbf{plantear} — \textquotedblleft{}to raise, to bring up; to pose\textquotedblright{}.

\begin{quote}
Quiero plantear una pregunta.
\end{quote}

\begin{quote}
\emph{I want to raise a question.}
\end{quote}

\emph{Plantear un problema} — to raise a problem. \emph{Plantearse} — to think about doing.

\begin{cousinweb}
\emph{plantear} is built on \emph{planta} 'sole of the foot; ground plan', a learned borrowing from Latin \emph{planta} (EtymDB): to set something down so all can see it.
\end{cousinweb}

\begin{grammarlens}[title={What you've built}]
One more word for telling, explaining and describing at B1.
\end{grammarlens}

\subsection*{Your turn}
\begin{itemize}
\raggedright
  \item \emph{Say it:} \emph{plantear}
  \item \emph{Say it:} \emph{plantear}, once more
  \item \emph{Say it:} a sentence of your own with \emph{plantear}
  \item \emph{Recall:} say the Spanish for the obstacle, then the Spanish for complicated, then say \emph{plantear} again
\end{itemize}

\begin{tcolorbox}[breakable,colback=teal!4,colframe=teal!35!black,title={Before you move on}]
What does plantear mean? (\textquotedblleft{}To raise, to bring up; to pose\textquotedblright{}.) Say it once more.
\end{tcolorbox}
